% 330_Operationen_T4_Abgrenzung_De_ch.tex
% Dok. 330, FFGFT-Korpus, August 2026 — Kapitelfile
% Eingebunden von Sources/wr_standalone_A4/330_Operationen_T4_Abgrenzung_De.tex

\providecommand{\K}{}\renewcommand{\K}{\textbf{[K]}}
\providecommand{\B}{}\renewcommand{\B}{\textbf{[B]}}
\renewcommand{\S}{\textbf{[S]}}
\providecommand{\SETZ}{}\renewcommand{\SETZ}{\textbf{[SETZUNG]}}
\providecommand{\E}{}\renewcommand{\E}{\textbf{[E]}}
\providecommand{\lP}{}\renewcommand{\lP}{\ell_{\mathrm{P}}}
\providecommand{\Tmass}{}\renewcommand{\Tmass}{\tilde{T}}

\definecolor{ffgftblue}{RGB}{30,90,160}
\definecolor{warncolor}{RGB}{140,0,40}
\tcbset{
  base/.style={breakable,enhanced,boxrule=0.6pt,arc=3pt,
    fonttitle=\bfseries\small,coltitle=white,top=4pt,bottom=4pt,left=6pt,right=6pt},
  ffgftbox/.style={base,colback=blue!5,colframe=ffgftblue,colbacktitle=ffgftblue},
  warnbox/.style={base,colback=red!5,colframe=warncolor,colbacktitle=warncolor}
}

\section*{Zusammenfassung}

Dieses Dokument konsolidiert die Abgrenzung der FFGFT gegenüber
diskreten Emergenzmodellen (hier: HLV nach Krüger, August 2026)
und klärt, welche Operationen in FFGFT zwischen $T^4$ und dem
Beobachtraum $\mathbb{R}^3 \times \mathbb{R}_t$ vermitteln.

Zentrale Aussage: Das Scheitern diskreter Gittermodelle (Klasse-X4-
Überlappungen, fehlende Operator-Spezifität, offener Kontinuums-Limes)
ist keine Überraschung, sondern eine strukturelle Konsequenz davon,
dass diese Modelle geometrische Hilfskonstrukte für ontologische
Bausteine halten. FFGFT vermeidet diese Fehler nicht durch eine
andere Wahl des Gitters, sondern dadurch, dass die Fundamentalebene
$T^4$ kompakt und geometrisch vollständig ist — und die drei Operationen,
die von ihr in den Beobachtraum führen, streng verschiedene logische
Status haben.

Die drei Operationen (Dok.~270 \K):
\begin{enumerate}
  \item \textbf{Typ~I} — Dualitäts-Dekompaktifizierung:
    $S^1_m \to \mathbb{R}_t$ via $\Tmass\cdot m = 1$.
    Informationserhaltend (Einbettung, keine Projektion).
  \item \textbf{Typ~II} — Geometrische Projektion:
    $D3 \to D2 \to D1$. Verlustbehaftet, nicht reversibel.
    Der Schritt $D3 \to D2$ ist als holographische Projektion
    identifiziert.
  \item \textbf{Typ~III} — Repräsentationale Übersetzung:
    $T^4 \leftrightarrow \mathcal{H}$ (Hilbertraum/Matrix).
    Bijektiv, verlustfrei — Aufwärtsbewegung, keine geometrische
    Dimensionsreduktion.
\end{enumerate}

% ============================================================
\section{Ausgangslage: Das HLV-Audit als Fallstudie}
\label{sec:hlv}

Marcel Krüger hat in \emph{Local-to-Global Compatibility on Aperiodic
Cut-and-Project Carriers} (v1.4, August 2026) das eigene Helix-Light-Vortex-Modell
(HLV) vier vorab eingefrorenen mathematischen Tests unterzogen.
Alle wesentlichen Behauptungen scheitern oder bleiben offen.

\begin{enumerate}
  \item \textbf{Lokal-zu-Global-Kompatibilität:}
    5.455 echte Volumen-Überlappungen (Klasse~X4) bei 3.384 projizierten
    Tetraedern. Sechs Zeugen mit 120-stelliger Arithmetik zertifiziert.
    \emph{Scheitern.} \E
  \item \textbf{Operator-Spezifität:}
    Ein gewöhnliches periodisches kubisches Gitter (CUBIC) übertrifft
    HLV im Spektral-Score (GAE-001A). \emph{Abgelehnt.} \E
  \item \textbf{Geometrie-zu-Konstitutivgesetz:}
    Die reine Trägergeometrie legt keine physikalischen Dichten,
    Massen oder Kopplungen fest (Proposition~2.2).
    \emph{Trivial negativ.} \E
  \item \textbf{Endlich-zu-Kontinuum:}
    Ein beschränkter Graph besitzt beschränktes Spektrum, keinen
    UV-Limes. Mosco-/Resolventen-Konvergenz unbewiesen. \emph{Offen.} \E
\end{enumerate}

\begin{tcolorbox}[warnbox,title={Korollar (Krüger 4.5) \E}]
  Lokale Gültigkeit eines Operators impliziert keine globale
  Einbettbarkeit des zugrundeliegenden Trägers.
\end{tcolorbox}

\subsection{Die strukturelle Ursache}

HLV versucht, geometrische Körper (3D-Tetraeder aus $\mathbb{R}^6$-Projektion)
als ontologische Bausteine des physikalischen Raums zu verwenden.
Damit erbt es alle vier Probleme:
Geometrische Elemente überlappen im Raum (1),
das Spektrum ist nicht gitterspezifisch (2),
aus der Geometrie allein folgt keine Physik (3),
und ein festes Gitter erzeugt keinen Kontinuumslimes (4).

Das Gegenmittel ist nicht ein \emph{besseres} Gitter.
Es ist ein fundamentaler Träger, der von vornherein
geometrisch vollständig und ohne Einbettungsproblem ist.

% ============================================================
\section{Die FFGFT-Fundamentalebene: $T^4$ mit D4-Gitter}
\label{sec:t4}

\subsection{Struktur von $T^4$}

Die Fundamentalebene der FFGFT ist der kompakte 4-Torus
\[
  T^4 = S^1_x \times S^1_y \times S^1_z \times S^1_m
\]
mit drei räumlichen Kreisen $S^1_{x,y,z}$ (Radius $R_H$, quasi-kontinuierlich)
und einem Masse-Kreis $S^1_m$ (Radius $R_m \sim \xi\cdot\lP$, diskret).
Die vierte Richtung $S^1_m$ ist nicht eine Massedimension neben einer
separaten Zeit, sondern die Ersatz-Zeit-Richtung selbst.
Masse und Zeit sind zwei Lesarten \emph{derselben} Dimension;
$\Tmass\cdot m = 1$ ist das Übersetzungswörterbuch (Dok.~270 \K,
Dok.~314 \K). \SETZ

Das Gitter auf $T^4$ ist das D4-Gitter (Dok.~314 \K):
Kissenzahl~24, Packungsdichte $|\text{Aut}(D_4)| = 1152$,
Trialität als eingebaute $\mathbb{Z}_3$-Symmetrie (16 Ordnung-3-Elemente
erzeugen $T^4/\mathbb{Z}_3$-Orbifold mit 9 Fixpunkten).
Ein Negativbefund mit Satzcharakter: $5 \nmid 1152$, also enthält
$\text{Aut}(D_4)$ kein Element der Ordnung~5 \K.

\subsection{Warum $T^4$ flach ist — intrinsische vs.\ extrinsische Krümmung \K}
\label{sec:flach}

Ein Torus trägt im Allgemeinen \emph{extrinsische} Krümmung, wenn er in
einen höherdimensionalen Umgebungsraum eingebettet wird (z.\,B.\
$T^2 \hookrightarrow \mathbb{R}^3$): die Außenseite des Rings ist
konvex (positive Gauß-Krümmung), die Innenseite konkav (negative
Gauß-Krümmung). Die \emph{intrinsische} Gauß-Krümmung hingegen, die das
Theorema Egregium als isometrische Invariante auszeichnet, addiert sich
punktweise zu null — positive und negative Krümmungsbeiträge heben sich
in der Summe auf.

\begin{tcolorbox}[ffgftbox,title={Intrinsisch flacher Torus \K}]
  $T^4$ ist intrinsisch flach: Ein Bewohner im Inneren misst keine
  Krümmung. Winkelsummen sind $\pi$, parallele Geodäten bleiben
  parallel. $T^4$ ist lokal isometrisch zu $\mathbb{R}^4$ —
  mit globaler Identifikation der Ränder, aber ohne lokale
  Krümmungskorrekturen.
\end{tcolorbox}

Dies hat zwei direkte Konsequenzen für die FFGFT:

\textbf{(a) D4-Gitter durch Topologie, nicht Krümmung.}
Die Auszeichnung des D4-Gitters gegenüber $\mathbb{Z}^4$ ist
\emph{nicht} durch Krümmungsunterschiede bedingt — der Torus ist
ohnehin flach (Dok.~314 \K: „Der Torus ist ohnehin flach —
die Krümmung war nie der Unterschied").
Was ins Spektrum wandert, sind die topologischen Daten:
Windungszahlen, Modeninhalt, Trialität. Die Unterscheidung
$\mathbb{Z}^4$/$D_4$ ist eine spektrale, keine geometrische.

\textbf{(b) Fourier-Analyse ohne Krümmungskorrekturen.}
Auf dem flachen Torus sind die Fourier-Moden $e^{ik \cdot x}$
die natürliche Basis; die Eigenwerte des Laplacians sind $|k|^2$
ohne Krümmungskorrektur. Dies ist die Grundlage dafür, dass die
Typ-III-Übersetzung $T^4 \leftrightarrow L^2(T^4) \otimes \mathbb{C}^3$
\emph{bijektiv und verlustfrei} ist: Die Fourier-Zerlegung auf einem
flachen Torus ist eine vollständige, isometrische Abbildung.
Ein gekrümmter Träger hätte Korrekturterme, die die Bijektivität
stören würden.

Dies unterscheidet FFGFT strukturell von Ansätzen wie HLV, die
geometrische Körper in $\mathbb{R}^3$ einbetten und dabei
Überlappungsprobleme erzeugen: Die Einbettungsfrage entsteht
in FFGFT gar nicht, weil $T^4$ als intrinsisch flacher,
in sich geschlossener Träger keine Einbettung in einen
Außenraum benötigt.



$T^4$ ist kompakt, vollständig und geometrisch in sich geschlossen.
Es gibt kein Einbettungsproblem, weil $T^4$ nicht in einem äußeren
Raum \emph{eingebettet} werden muss: Es \emph{ist} die Fundamentalebene.
Die Klasse-X4-Fehler von HLV entstehen, weil dort Projektionselemente
in einen externen $\mathbb{R}^3$ eingebettet werden sollen.
In FFGFT stellt sich diese Frage nicht.

\subsection{Typ~III: Repräsentationale Übersetzung nach oben}

Die Übersetzung $T^4 \leftrightarrow \mathcal{H} = L^2(T^4) \otimes \mathbb{C}^3$
(Dok.~230, 314, 322) ist eine bijektive, verlustfreie Aufwärtsbewegung
(Typ~III, Dok.~270 \K).
Sie fügt keine geometrischen Extra-Dimensionen hinzu, sondern drückt
denselben $T^4$-Inhalt in einer höherdimensionalen Sprache aus.
Das Spektrum des fundamentalen fraktalen Operators $\hat{F}$
(Dok.~322, selbstadjungiert [B], Dok.~327) hat die Teilchenmassen
als Eigenwerte; $\xi = \lambda_{\min}(\hat{F}_{D_4})$ ist der kleinste
Eigenwert \B.

Die Gitterstruktur geht dabei nicht verloren, sondern wandert
ins Spektrum und in die Faser \K:
Das D4-Spektrum trennt, was die Krümmung nicht trennt
($\mathbb{Z}^4$ vs.\ $D_4$ im Modeninhalt);
Trialität ist eingebaut, nicht angeheftet.

% ============================================================
\section{Die drei Operationen von $T^4$ zum Beobachtraum}
\label{sec:operationen}

Dok.~270 unterscheidet streng drei Operationstypen, die zusammen
die beobachtbare Physik aus $T^4$ ableiten:

\subsection{Typ~I: Dualitäts-Dekompaktifizierung}

\[
  S^1_m \;\xrightarrow{\;\Tmass\cdot m = 1\;}\; \mathbb{R}_t
\]

Das Abrollen der kompakten Massedimension in die nichtkompakte
Zeitrichtung ist eine \emph{Einbettung}, keine Projektion.
Kein Modeninhalt geht verloren; die Zeit entsteht als zusätzliche
nichtkompakte Richtung.
Was nicht sichtbar bleibt, ist allein die Kompaktheit — sie erscheint
indirekt als Diskretheit des Spektrums, kodiert durch $\xi^n$ \K.

Dies beantwortet einen Einwand aus Krügers Punkt~4:
Das Kontinuumsproblem entsteht, wenn man versucht, einen
\emph{festen endlichen Graphen} in ein Kontinuum übergehen zu lassen.
In FFGFT ist die Zeitrichtung kein Limes eines Graphen,
sondern das Ergebnis einer Einbettung ($S^1_m \hookrightarrow \mathbb{R}_t$),
die von Anfang an stetig ist.

\subsection{Typ~II: Geometrische Projektion (verlustbehaftet)}

Die weiteren Schritte $D3 \to D2 \to D1$ sind echte geometrische
Projektionen räumlicher Richtungen: verlustbehaftet, nicht reversibel.
Der Schritt $D3 \to D2$ wird als holographische Projektion
identifiziert (Dok.~270 \K).

\begin{tcolorbox}[ffgftbox,title={Einordnung des holographischen Schritts}]
  Der Schritt $D3 \to D2$ ist \emph{nicht} die Fundamentalebene von FFGFT.
  Er ist eine verlustbehaftete Beobachterprojektion auf einer
  Unterebene — physikalisch relevant für Oberflächenphänomene
  (Schwarze-Loch-Membran, Dok.~325), aber nicht für die
  Herleitung von Massen und Kopplungen, die auf $T^4$ verankert sind.
  Die zweidimensionale Projektionsfläche ist kein fundamentales Objekt.
\end{tcolorbox}

\subsection{Gegenüberstellung der drei Typen}

\begin{table}[h!]
\small
\begin{tabular}{@{}llll@{}}
\toprule
Typ & Operation & Richtung & Status \\
\midrule
I   & Dualitäts-Dekompaktifizierung $S^1_m \to \mathbb{R}_t$
    & geometrisch abwärts & informationserhaltend \K \\
II  & Geometrische Projektion $D3 \to D2 \to D1$
    & geometrisch abwärts & verlustbehaftet \K \\
III & Repräsentationale Übersetzung $T^4 \leftrightarrow \mathcal{H}$
    & repräsentational aufwärts & bijektiv, verlustfrei \K \\
\bottomrule
\end{tabular}
\end{table}

Die scheinbare Spannung — Projektion verliert Information (A),
Kompaktifizierung dekompaktifiziert (B), Aufwärtsübersetzung ist
verlustfrei (C) — löst sich durch die Typentrennung auf (Dok.~270).

% ============================================================
\section{Warum Krügers vier Probleme in FFGFT nicht auftreten}
\label{sec:vergleich}

\begin{enumerate}
  \item \textbf{Klasse-X4-Überlappungen:}
    Entstehen, weil HLV Projektionselemente in $\mathbb{R}^3$ einbetten
    will. In FFGFT ist $T^4$ die Fundamentalebene; es gibt keine
    Einbettung in einen externen Raum.

  \item \textbf{Fehlende Operator-Spezifität:}
    In FFGFT ist das Spektrum nicht durch ein beliebiges Gitter,
    sondern durch das D4-Gitter mit eingebauter Trialität bestimmt.
    Der Negativbefund $5 \nmid |\text{Aut}(D_4)|$ ist selbst ein
    spezifischer Satz über das Spektrum \K.

  \item \textbf{Geometrie ohne Physik:}
    $T^4$ allein legt ebenfalls keine Kopplungskonstanten fest.
    Die Physik kommt aus der Kombination $T^4$-Spektrum $+$
    $\xi = \lambda_{\min}(\hat{F}_{D_4})$ als Selektionsprinzip.
    Das ist kein Widerspruch zu Krügers Punkt~3, sondern seine
    Auflösung: nicht die Geometrie allein, sondern der
    minimale Eigenwert des Operators auf der Geometrie bestimmt die Physik.

  \item \textbf{Kontinuumslimes:}
    Kein Limes-Problem, weil Typ~I eine Einbettung ist (kein
    Grenzübergang nötig) und Typ~III bijektiv ist (kein Informationsverlust).
    Einzig Typ~II ist verlustbehaftet, aber das ist bekannt und
    im Korpus ausgewiesen (Dok.~270).
\end{enumerate}

\subsection{Wie $\Tmass\cdot m = 1$ im Hilbertraum erscheint \K}
\label{sec:dualitaet-hilbert}

Die Zeit-Masse-Dualität $\Tmass\cdot m = 1$ erscheint im Hilbertraum
$\mathcal{H} = L^2(T^4)\otimes\mathbb{C}^3$ nicht als Zeitentwicklung,
sondern in drei komplementären Formen:

\textbf{(a) Als vierte Wicklungszahl $n_4$.}
Der Multiindex $n = (n_1,n_2,n_3,n_4)\in\mathbb{Z}^4$ zählt die
Wicklungen in allen vier Torusrichtungen \B\ (Dok.~322).
Die Massedimension $S^1_m$ trägt die Wicklungszahl $n_4$.
Die Reziprozitätsrelation $\lambda_4\cdot m = 2\pi$ gilt modenweise
exakt \K\ (Dok.~314): Wicklungszahl $=$ Massenzahl.
$\Tmass\cdot m = 1$ ist damit die Kaluza-Klein-Relation der
vierten Torusrichtung, ablesbar als $n_4$-Quantenzahl im Spektrum.

\textbf{(b) Als Massenoperator-Eigenwertgleichung.}
\begin{equation}
  \hat{M}\,|n_\theta,n_\phi,p\rangle = m_i\,|n_\theta,n_\phi,p\rangle
  \quad\textbf{[K]}
\end{equation}
Die Eigenwerte sind die Teilchenmassen (Dok.~322, Dok.~320).
Die Dualität erscheint als Spektrum des Operators — nicht als
separate Gleichung neben dem Hilbertraum, sondern als seine
Eigenwertstruktur.

\textbf{(c) Als Radienklassentrennung im Spektrum.}
Die Aufspaltung bei $R_4 \neq R_{123}$ macht $\Tmass\cdot m = 1$
direkt ablesbar \K\ (Dok.~314): Die Massedimension erzeugt
einen großen Modenabstand $\sim 1/R_m$, die Raumrichtungen
einen quasi-kontinuierlichen Abstand $\sim 1/R_H$.
Das ist die strukturelle Trennung von Masse und Raum im Hilbertraum.

\begin{tcolorbox}[ffgftbox,title={Was im Hilbertraum \emph{nicht} erscheint}]
  Die Zeit $\mathbb{R}_t$ selbst. Sie ist das Ergebnis der
  Typ-I-Dekompaktifizierung $S^1_m \to \mathbb{R}_t$ — einer
  geometrischen Abwärtsoperation \emph{vor} der Typ-III-Übersetzung.
  Im Hilbertraum $L^2(T^4)$ lebt nur der Torus;
  die Zeit tritt erst auf, wenn man aus dem Hilbertraum herausgeht
  und das Spektrum in $\mathbb{R}_t$ einbettet.
  $\Tmass\cdot m = 1$ steckt im Hilbertraum als $n_4$-Wicklungszahl
  und Massenoperator-Eigenwertstruktur, nicht als Zeitentwicklung.
\end{tcolorbox}

\subsection{Wo die fraktale Eigenschaft im Hilbertraum steckt \K\ / \B}
\label{sec:fraktal-hilbert}

Die fraktale Eigenschaft der FFGFT erscheint im Hilbertraum an
drei präzise trennbaren Stellen:

\textbf{(a) Im Operator $\hat{F}$ selbst — als Selbstähnlichkeitsstruktur \B.}
\begin{equation}
  \hat{F} = \bigoplus_{n=1}^{N} S_n\circ L_n,\quad
  S_n = r_n U_n,\quad r_n = \xi^n \in (0,1)
\end{equation}
Die Skalenoperatoren $S_n$ implementieren diskrete Skalentransformationen;
jede Stufe $n$ ist eine Kontraktion um $r_n$ — das ist
Selbstähnlichkeit als Operator.
$\hat{F}$ kodiert die Sub-Planck-Geometrie als hierarchische Faltung
von Skalenebenen bis zur Mindestlängenskala $L_0 = \xi\cdot\lP$
(Dok.~180, 327 \K).

\textbf{(b) In der fraktalen Dimension $D_f = 3-\xi$ — als Spektraleigenschaft \K.}
\begin{equation}
  D_f = \inf\!\left\{s:\sum_i\lambda_i^s < \infty\right\}
  = 3 - \xi \approx 2{,}9999
  \quad\textbf{[K]}
\end{equation}
$D_f$ ist über das Hausdorff-Maß des Spektrums von $\hat{F}$
definiert (Dok.~322): das Spektrum auf $L^2(T^4)$ hat eine fraktale
Feinstruktur. Der Wert ist räumlich ($3-\xi$, drei Dimensionen),
nicht vierdimensional.
Lokale Wirkung: $D_f^\text{Raum} = 3-\xi$ mit $6{,}7\times10^{-5}$
relativ (Dok.~314 \K).

\textbf{(c) Im Interface-Faktor $K_\text{frak}$ — als Skalenbrücke \K.}
\begin{equation}
  K_\text{frak} = 1 - 100\xi = \tfrac{74}{75} \approx 0{,}9867
\end{equation}
$K_\text{frak}$ ist kumulativ (Dok.~133): es akkumuliert über den
RG-Lauf von 17 Größenordnungen und wirkt beim Übergang von der
geometrischen zur gemessenen (SI-)Masse:
$m_i^\text{SI} = K_\text{frak}\cdot C_\text{conv}\cdot m_i^\text{bare}$.

\begin{tcolorbox}[ffgftbox,title={Entscheidende Trennung \K}]
  Im \emph{eingerollten} $T^4$ (Hilbertraum) gibt es
  \textbf{keine} fraktalen Korrekturen —
  Wicklungszahlen sind rein topologisch (Dok.~314 \K).
  Die fraktale Eigenschaft tritt erst beim Ausrollen
  (Typ~I: $S^1_m \to \mathbb{R}_t$) als akkumulierter Defekt auf.
  Im Hilbertraum selbst ist $\hat{F}$ ein beschränkter
  selbstadjungierter Operator mit reellem Spektrum \B\ —
  die Fraktalität steckt in der Skalenhierarchie der Operatorstruktur
  ($r_n = \xi^n$), nicht im Träger $L^2(T^4)$.
\end{tcolorbox}


\label{sec:hilbert}

Der flache Torus $T^4$ erlaubt eine saubere Fourier-Analyse
ohne Krümmungskorrekturen. Drei Fragen stellten sich,
ob die Struktur des Hilbertraums $\mathcal{H} = L^2(T^4)\otimes\mathbb{C}^3$
mit D4-Gitter, $\mathbb{Z}_3$-Orbifold und fraktalem Operator $\hat{F}$
konsistent ist. Alle drei sind durch den Korpus und zwei
Prüfskripte (330-1, 330-2) entschieden.

\subsection{Frage 1: Haar-Maß und D4-Gitter \B}

Das kanonische Haar-Maß auf $T^4$ (translationsinvariantes
Lebesgue-Maß) ist mit der D4-Spektralstruktur kompatibel:
\begin{itemize}
  \item Fourier-Moden sind unter dem Haar-Integral exakt orthogonal
    (Skript 330-1, A1).
  \item Das D4-Spektrum enthält ausschließlich gerade Normschalen
    $|k|^2 \in \{2,4,6,8,\ldots\}$ — ungerade fehlen vollständig
    (A2, Dok.~314 \K).
  \item $\text{Aut}(D_4) \subset \text{Iso}(T^4)$:
    das Haar-Maß ist unter D4-Automorphismen invariant (A3).
  \item Der $\mathbb{Z}_3$-Orbifold bricht die Translationsinvarianz
    und spaltet die 24er-Entartung zu $9\oplus8\oplus4\oplus2\oplus1$
    auf (A4, Dok.~314 \K) — aber nur als Feinstruktur
    $\sim\xi \approx 10^{-4}$, unterhalb jeder Messschwelle (A5).
\end{itemize}
Die Container-Struktur $4\times6$ bleibt praktisch intakt. \B

\subsection{Frage 2: Fixpunkte des Orbifolds \K\ / \S}

Die neun Fixpunkte von $T^4/\mathbb{Z}_3$ folgen aus der
Lefschetz-Fixpunktformel (Skript 330-2, B1):
\[
  |\det(1-A)| = |(1-\omega)(1-\omega^2)|^2 = 3^2 = 9
\]
für die $\mathbb{Z}_3$-Wirkung mit Eigenwerten
$\{\omega,\omega^2,\omega,\omega^2\}$ auf $\mathbb{C}^2$.
Die drei neutrinorelevanten Fixpunkte sind (Dok.~322 \K):
\[
  \nu_1 \leftrightarrow F_2,\quad
  \nu_2 \leftrightarrow F_5,\quad
  \nu_3 \leftrightarrow F_7.
\]
$F_7$ unterscheidet sich durch $n_4 \neq 0$ (Masserichtung)
von $F_2,F_5$; die drei Fixpunkte sind paarweise verschieden (B2).

Was bleibt \S: Die Rückreaktion der Hawking-Emission auf die
Fixpunktstruktur während $M \to M_\mathrm{coll} = \mP/\sqrt{2}$
ist in Dok.~325 explizit als offen deklariert.
Insbesondere: Wie verändert sich die Neutrino-Lokalisierung,
wenn der Horizont auf $r_s = \sqrt{2}\,\lP$ schrumpft?

\subsection{Frage 3: Fraktale Gewichtung $r_n$ und $L^2(T^4)$ \B}

Der fraktale Operator $\hat{F} = \bigoplus_{n=1}^N S_n\circ L_n$
mit $S_n = r_n U_n$, $r_n \in (0,1)$ (Skript 330-2, B4--B5):
\begin{itemize}
  \item $r_n = \xi^n$: analytisch positiv für alle endlichen $n$
    (log $r_n = n\ln\xi < 0$), $N$ endlich durch $L_0 = \xi\cdot\lP$
    (Dok.~180, 327 \K).
  \item $\|\hat{F}\| \leq \sum r_n < \infty$: $\hat{F}$ ist beschränkt
    und damit auf $\mathcal{H} = L^2(T^4)\otimes\mathbb{C}^3$
    selbstadjungiert \B\ (Dok.~327).
  \item $r_n = f(n)$ ist \emph{positionsunabhängig}: $\hat{F}$ bricht
    die Translationsinvarianz von $L^2(T^4)$ \emph{nicht}.
    Die Translationsbrechung kommt ausschließlich vom $\mathbb{Z}_3$-Orbifold
    (Frage 1, A4) — und zwar nur als Feinstruktur $\sim\xi$.
  \item Die $\mathbb{Z}_3$-Paarung $(k,-k)$ erzwingt $\hat{F} = \hat{F}^\dagger$
    auf dem flachen $L^2(T^4)$ \B\ (Dok.~327, Defektindizes $(0,0)$).
\end{itemize}

\begin{tcolorbox}[ffgftbox,title={Statusbilanz der drei Fragen}]
  Frage 1 (Haar/D4): \B\ — vollständig gelöst.\\
  Frage 2 (Fixpunkte): \K\ für 9 Fixpunkte und $\nu$-Lokalisierung;
  \S\ für Rückreaktion bei Verdampfung (explizit offen, Dok.~325).\\
  Frage 3 ($r_n$/$\hat{F}$): \B\ — $\hat{F}$ beschränkt und selbstadjungiert;
  $r_n$ bricht Translationsinvarianz strukturell nicht.
\end{tcolorbox}

% ============================================================
\section{Was aus der Gegenüberstellung mit diskreten Emergenzmodellen bleibt}
\label{sec:bilanz}

Das HLV-Audit (Krüger 2026) und die Klärung der Hilbertraum-Fragen
ergeben zusammen eine klare Bilanz:

\textbf{Strukturell gesichert:} Die vier HLV-Probleme
(Überlappung, Spezifität, Konstitutivgesetz, Kontinuumslimes)
treten in FFGFT nicht auf — wegen $T^4$ als intrinsisch flachem,
in sich geschlossenem Träger (kein Einbettungsproblem), der
bijektiven Typ-III-Übersetzung (kein Kontinuumslimes nötig)
und der Typ-I-Einbettung $S^1_m \hookrightarrow \mathbb{R}_t$
(kein Graphen-zu-Kontinuum-Problem).

\textbf{Korrekt aus der vergleichenden Diskussion:}
Die holographische Projektion $D3\to D2$ (Typ~II) ist verlustbehaftet
und nicht fundamental; Masse und Zeit sind nicht in der 2D-Fläche
kodiert, sondern entstehen durch Typ-I ($\Tmass\cdot m = 1$) und
die Rekursion $\xi^n$ — die Fläche ist Beobachterprojektion, kein Substrat.

\textbf{Korrektur gegenüber unkritischer Übernahme:}
Es gibt in FFGFT keine „zweidimensionale Projektionsfläche als
fundamentales Objekt". Der Korpus hat eine klare ontologische
Hierarchie: $T^4$ als primäre Realität (Dok.~152 \SETZ),
nicht eine 2D-Fläche. $\Tmass\cdot m = 1$ ist die
Kaluza-Klein-Relation der Massedimension, kein
„Leseprinzip auf einem Schirm".

\vspace{4pt}\noindent
\textbf{Prüfskripte:}
\texttt{pruef\_330\_1\_haar\_d4.py} (7/7),
\texttt{pruef\_330\_2\_fixpunkte\_rn.py} (9/9).



\begin{table}[h!]
\small
\begin{tabular}{@{}lc@{}}
\toprule
Element & Status \\
\midrule
$T^4 = S^1_x \times S^1_y \times S^1_z \times S^1_m$ als Fundamentalebene & \SETZ \\
D4-Gitter auf $T^4$; Kissenzahl~24; $|\text{Aut}(D_4)| = 1152$ & \K \\
$\Tmass\cdot m = 1$ als KK-Relation der Massedimension & \K \\
Drei Operationstypen (Dok.~270): Typ~I/II/III & \K \\
Typ~I informationserhaltend (Einbettung) & \K \\
Typ~III bijektiv ($T^4 \leftrightarrow \mathcal{H}$, Dok.~322/314) & \K \\
Typ~II verlustbehaftet; $D3\to D2$ holographisch & \K \\
$5 \nmid |\text{Aut}(D_4)|$: kein Ordnung-5-Element & \K \\
HLV-Klasse-X4-Fehler strukturell ausgeschlossen (kein Einbettungsproblem) & \B \\
Kontinuumslimes strukturell gelöst durch Typ~I (Einbettung) & \K \\
\bottomrule
\end{tabular}
\end{table}

\vspace{6pt}\noindent\small
\textcopyright\ 2026 Johann Pascher · \href{https://creativecommons.org/licenses/by/4.0/}{CC BY 4.0}
